%% version for arxiv submission 31/8/2012
%% copy of kpz-22.tex

\documentclass[11pt,a4paper,draft,reqno]{amsart}
%\usepackage{pdfsync}
\usepackage{geometry} % see geometry.pdf on how to lay out the page. There's lots.
%\geometry{a4paper} % or letter or a5paper or ... etc
%\geometry{margins=.5in} % for example, change the margins to 2 inches all round
%\geometry{landscape} % set up the page for landscape
% read geometry.pdf for detailed page layout information
\usepackage{amsmath,amsthm,amssymb}
%\usepackage{showkeys}
\usepackage{amsrefs}
\usepackage{hyperref}
\usepackage[latin1]{inputenc}
\usepackage{color}

%\usepackage[firstpage]{draftwatermark}
%\SetWatermarkScale{5}

\newtheorem{definition}{Definition}
\newtheorem{lemma}{Lemma}
\newtheorem{corollary}{Corollary}
\newtheorem{theorem}{Theorem}

\newcommand{\TT}{\mathbb{T}}
\newcommand{\PP}{\mathbb{P}}
\newcommand{\II}{\mathbb{I}}
\newcommand{\RR}{\mathbb{R}}
\newcommand{\ZZ}{\mathbb{Z}}
\newcommand{\CC}{\mathbb{C}}
\newcommand{\expect}{\mathbb{E}}
\newcommand{\cL}{\mathcal{L}}
\newcommand{\cE}{\mathcal{E}}
\newcommand{\cF}{\mathcal{F}}
\newcommand{\kpz}{\mathrm{kpz}}
\newcommand{\she}{\mathrm{she}}
\newcommand{\dd}{\mathrm{d}}
\newcommand{\Lz}{{L^2_0(\TT)}}
\newcommand{\cS}{\mathcal{S}}
\newcommand{\cN}{\mathcal{N}}
\newcommand{\cA}{\mathcal{A}}
\newcommand{\cQ}{\mathcal{Q}}
\newcommand{\cC}{\mathcal{C}}
\newcommand{\cR}{\mathcal{R}}
\newcommand{\cM}{\mathcal{M}}
\newcommand{\eps}{\varepsilon}
\newcommand{\cyl}{\mathcal{C}yl}
\newcommand{\la}{\langle}
\newcommand{\ra}{\rangle}

\newcommand{\Lou}{L^{\textrm{ou}}}
\newcommand{\tr}{\mathrm{Tr}}

\newcommand{\Q}{\ZZ^2\backslash\{0\}}
\newcommand{\QN}{\ZZ^2_N\backslash\{0\}}


%%% BEGIN DOCUMENT
\begin{document}

\title[Regularization for SBE]{Regularization by noise and \\  stochastic Burgers equations}
\author{M.~Gubinelli}
\address[M.~Gubinelli]{CEREMADE  UMR 7534 -- Universit\'e Paris--Dauphine}
\email[M.~Gubinelli]{massimiliano.gubinelli@ceremade.dauphine.fr}
\author{M.~Jara}
\address[M.~Jara]{IMPA\\
Estrada Dona Castorina 110\\
CEP 22460-320\\
Rio de Janeiro\\
Brazil}
\email[M.~Jara]{mjara@impa.br}
%\date{v17 20120531}
%\dateposted{XXX}

\begin{abstract}
For any $\theta > 1/2$ we study a generalized 1d periodic SPDE of Burgers type:
$$
\partial_t u_t =- A^\theta u_t  + F(u_t)  + A^{\theta/2} \partial_t W_t 
$$
where $-A$ is the 1d Laplacian,  $F(u_t)(\xi)=\partial_\xi (u_t(\xi))^2$ is the Burgers nonlinearity, $\partial_t W_t$ is a space-time white noise and  the initial condition $u_0$ is taken to be (space) white noise. We propose a notion of solution for this equation in the stationary setting. For these solutions we point out how the noise provide a regularizing effect allowing to prove existence and suitable estimates when $\theta>1/2$. When $\theta>5/4$ we obtain easily uniqueness. 
\end{abstract}
\keywords{Kardar--Parisi--Zhang equation, SPDEs, noise regularization}
\subjclass[2000]{00X00}

\maketitle

%\section*{Introduction}
Consider the stochastic Burgers equation (SBE) on the one dimensional torus $\TT=(-\pi,\pi]$
\begin{equation}
\label{eq:burgers}
\dd u_t = \frac12 \partial_\xi^2 u_t(\xi) \dd t + \frac12 \partial_\xi (u_t(\xi))^2 \dd t + \partial_\xi \dd W_t 
\end{equation}
where $W_t$ is a cylindrical white noise of the form
$
W_t(\xi) = \sum_{k\in\ZZ_0} e_k(\xi) \beta^k_t
$
with $\ZZ_0 = \ZZ\backslash \{0\}$ and $e_k(\xi)=e^{i k\xi}/\sqrt{2\pi}$ and $\{\beta_t^k\}_{t\ge 0, k\in\ZZ_0}$ is a family of complex Brownian motions such that $(\beta^k_t)^*=\beta^{-k}_t$ and with covariance $\expect[\beta_t^k \beta_t^q ]=\II_{q+k=0}$.   Formally the solution $u$ of eq.~\eqref{eq:burgers} is the derivative of the solution  of the Kardar--Parisi--Zhang equation
\begin{equation}
\label{eq:kpz}
\dd h_t = \frac12 \partial_\xi^2 h_t(\xi) \dd t + \frac12  (\partial_\xi h_t(\xi))^2 \dd t +  \dd W_t 
\end{equation}
which is believed to capture the macroscopic behavior of a large class of surface growth phenomena~\cite{KPZ}.

The main difficulty with eq.~\eqref{eq:burgers} is given by the rough nonlinearity which is incompatible with the distributional nature of the typical trajectories of the process. Note in fact that, at least formally, eq.~\eqref{eq:burgers} preserves the white noise on the Hilbert space $H=\Lz$ of square integrable, mean zero real function on $\TT$ and that the square in the non-linearity is almost surely $+\infty$ on the white noise. Additive renormalizations in the form of Wick products are not enough to cure this singularity~\cite{DDT}.

In~\cite{BG} Bertini and Giacomin showed that a particular regularization of~\eqref{eq:burgers}  converges in law to a limiting process $u^{\textrm{hc}}_t(\xi)=\partial_\xi \log Z_t(\xi)$ (which is referred to as the Hopf-Cole solution)  where $Z$ is the solution of the stochastic heat equation with multiplicative space--time white noise
\begin{equation}
\label{eq:she}
\dd Z_t = \frac12 \partial_\xi^2 Z_t(\xi) \dd t + Z_t(\xi) \dd W_t(\xi) .
\end{equation}

The Hopf--Cole solution is believed to be the correct physical solution for~\eqref{eq:burgers} however up to recently a rigorous notion of solution to~\eqref{eq:burgers} was lacking so the issue of uniqueness remained open. 

 Jara and Goncalves~\cite{JG} introduced a notion of \emph{energy solution} for eq.~\eqref{eq:burgers}  and showed that  the macroscopic current fluctuations of a large class of weakly non-reversible particle systems on $\ZZ$ obey the Burgers equation in this sense. Moreover their results show that also the Hopf-Cole solution is an energy solution of eq.~\eqref{eq:burgers}.

More recently Hairer~\cite{Hairer} obtained a complete existence and uniqueness result for KPZ. In this remarkable paper he uses the theory of controlled rough paths to give meaning to the nonlinearity and a careful analysis,  via systematic use of Feynman diagrams, of the series expansion of the solution. He was able to show that the solution he identifies  coincides  with the Cole-Hopf solution.

In this paper we take a different approach to the problem: we want to show that the stochastic part of the equation has a regularizing effect on the non-linear part. This is linked to some similar remarks of Assing~\cite{assing1,assing2}. Our point of view is motivated also by similar analysis in the PDE and SPDE context where the noise or a dispersive term provide enough regularization to treat some non-linear term: there are examples involving the stochastic transport equation~\cite{FGP}, the periodic Korteweg-de~Vries equation~\cites{kdv,babin-kdv} and the fast rotating Navier-Stokes equation~\cite{babin-ns}. In particular in the paper~\cite{kdv} it is shown how, in the context of the periodic Korteweg-de~Vries equation, an appropriate notion of controlled solution can make sense of the non-linear term in a space of distributions. This point of view has also links with the approach via controlled paths to the theory of rough paths~\cite{controlling}.


In our approach we are not able to obtain uniqueness for the SBE above and  we resort to study the more general equation (SBE$_\theta$):
\begin{equation}
\label{eq:burgers-theta}
\dd u_t = - A^\theta u_t \dd t + F(u_t) \dd t + A^{\theta/2} \dd W_t 
\end{equation}
where $F(u_t)(\xi)=\partial_\xi (u_t(\xi))^2$, $-A$ is the Lapacian with periodic b.c.{\color{red},} where $\theta\ge 0$ and where the initial condition is taken to be white noise. In the case $\theta=1$ we essentially recover the stationary case of the SBE above (modulo a mismatch in the noise term which do not affect its law). 

For any $\theta \ge 0$ we introduce a class $\mathcal{R}_\theta$ of distributional processes "controlled" by the noise, in the sense that these processes have a  \emph{small time} behaviour similar to that of the stationary  Ornstein-Uhlenbech process $X$  which solves the linear part of the dynamics:
\begin{equation}
\label{eq:ou-theta}
\dd X_t = - A^\theta X_t \dd t + A^{\theta/2} \dd W_t,
\end{equation}
where $X_0$ is white noise.
 When $\theta > 1/2$ we are able to show that the \emph{time integral} of the non-linear term appearing in SBE$_\theta$ is well defined, namely that for all $v\in \mathcal{R}_\theta$ 
$$
A^v_t = \int_0 ^t F(v_s) \dd s
$$
is a well defined process with continous paths in the space of distributions. Note that this process is not necessarily of finite variation with respect to the time parameter. 

A solution of the equation will then be naturally defined in this class and somehow coincides with the notion of energy solution introduced by Jara and Gon\c{c}alves. Solutions will then exists when $\theta > 1/2$. We are also able to show easily uniqueness when $\theta > 5/4$ but the  case $\theta=1$ seems still (way) out of range for this technique. 

Similar regularization phenomena for stochastic transport equations are studied in~\cite{FGP} and in~\cite{DF} for infinite dimensional SDEs. This is also linked to the fundamental paper of Kipnis and Varadhan~\cite{KV} on CLT for additive functionals and to the Lyons-Zheng representation for diffusions with singular drifts~\cites{MR1988703, MR2065168}.

\bigskip
\textbf{Plan.} In
Sec.~\ref{sec:controlled} we define the class of controlled paths and we recall some results of the stochastic calculus via regularization which are needed to handle the It\^o formula for the controlled processes.
Sec.~\ref{sec:ito-trick} is devoted to introduce our main tool which is a moment estimate of an additive functional of a stationary Dirichlet process in terms of the quadratic variation of suitable forward and backward martingales. In Sec.~\ref{sec:estimates} we use this estimate to provide uniform bounds for the drift of any stationary solution. These bounds are then used in Sec.~\ref{sec:existence} to prove tightness of the approximations when $\theta > 1/2$ and to characterize any accumulation point as a controlled solution of the stochastic Burgers equation. Finally in Sec.~\ref{sec:uniq} we prove our partial uniqueness result in the case $\theta > 5/4$. In Sec.~\ref{sec:alternative} we discuss related results for the model introduced in~\cite{DDT}.


\bigskip
\textbf{Notations.} We write $X \lesssim_{a,b,\dots} Y$ if there exists a positive constant $C$ depending only on $a,b,\dots$ such that $X \le C Y$. We write $X \sim_{a,b,\dots} Y$ iff $X\lesssim_{a,b,\dots} Y \lesssim_{a,b,\dots} X$. 

We let $\cS$ the space of smooth test functions on $\TT$, $\cS'$ the space of distributions and $\langle \cdot,\cdot\rangle$ the corresponding duality, the function  $\rho:\RR\to\RR$ is a positive smooth test function with unit integral and $\rho^\eps(\xi)=\rho(\xi/\eps)/\eps$ for all $\eps>0$.


On the Hilbert space $H=\Lz$ the family $\{e_k\}_{k\in\ZZ_0}$ is a complete orthonormal basis. On $H$ we consider the space of smooth cylinder functions  $\cyl$ which depends only on finitely many coordinates on the basis $\{e_k\}_{k\in \ZZ_0}$ and for $\varphi \in\cyl$ we consider the gradient $D \varphi : H\to H$ defined as $D \varphi(x) = \sum_{k\in\ZZ_0} D_k \varphi(x) e_k$ where $D_k = \partial_{x_k}$ and  $x_k = \langle e_k,x \rangle$ are the coordinates of $x$.
We let $A=-\partial_\xi^2$ and $B=\partial_\xi$ as  unbounded operators acting on  $H$. Note that $\{e_k\}_{k\in\ZZ_0}$ is a basis of eigenvectors of $A$ for which we denote $\{\lambda_k = |k|^2 \}_{k\in\ZZ_0}$ the associated eigenvalues. 
The linear operator $\Pi_N: H \to H$ is the projection on the subspace generated by $\{e_k\}_{k\in\ZZ_0, |k|\le N}$. 

For any $\alpha\in \RR$ define the space $\cF L^{p,\alpha}$ of functions on the torus for which
$$
|x|_{\cF L^{p,\alpha}} = \big[\sum_{k\in\ZZ_0} (|k|^\alpha |x_k|)^p\big]^{1/p}<+\infty
\, \text{ 
if $p<\infty$ and 
 }\, 
|x|_{\cF L^{\infty,\alpha}} = \sup_{k\in\ZZ_0} |k|^\alpha |x_k| <+\infty .
$$
Denote $\cC_T V = C([0,T],V)$ the space of continuous functions from $[0,T]$ to the Banach space $V$ endowed with the supremum norm and with $\cC^\gamma_T V = C^\gamma([0,T],V)$ the subspace of $\gamma$-H\"older continuous functions in $\cC_T V$ with the $\gamma$-H\"older norm. 



\section{Controlled processes}
\label{sec:controlled}


We introduce a space of stationary processes which "looks like" an Ornstein-Uhlenbeck process.

\begin{definition}[Controlled process]
\label{def:controlled}
For any $\theta\ge 0$ let $\mathcal{R}_\theta$ be the space of stationary stochastic processes $(u_t)_{0 \leq t \leq T}$ with  continuous paths in $\cS'$ such that 
\begin{itemize}
\item[i)] the law of $u_t$ is the white noise for all $t$;
\item[ii)] there exists a process $\cA \in C([0,T],\cS')$ of zero quadratic variation such that $\cA_0 = 0$ and satisfy the equation
$$
u_t(\varphi) = u_0(\varphi) + \int_0^t u_s(-A^\theta \varphi) \dd s+\cA_t(\varphi) + M_t(\varphi)
$$
for any test function $\varphi \in \mathcal S$, where $M_t(\varphi)$ is a martingale with quadratic variation $[M(\varphi)]_t =  2t\|A^{\theta/2} \varphi\|_{L^2_0(\TT)}^2$;
\item[iii)] the reversed processes $\hat u_t = u_{T-t}$, $\hat \cA_t = -\cA_{T-t}$ satisfy the same equation with respect to its own filtration (the backward filtration of $u$).
\end{itemize}
\end{definition}

For controlled processes we will prove that if $\theta>1/2$ the Burgers drift is well defined in the sense that, for any smooth test fuction $\varphi$,
$$
\lim_{\eps\to 0} \int_0^t \langle \varphi, F(\rho^\eps* u_s)\rangle \dd s
$$
exists in any $L^p(\PP)$. We denote with $\int_0^t  F( u_s) \dd s$ the resulting process  with values in space distributions. It will turn out that for this process we have a good control of its space and time regularity and also some exponential moment estimates.  Then it is relatively natural to \emph{define} solutions of eq.~\eqref{eq:burgers-theta} as follows:

\begin{definition}[Controlled solution]
A process $u\in\mathcal{R}_\theta$ is a \emph{controlled solution} of the generalized stochastic Burgers equation~\eqref{eq:burgers-theta} if
$$
 \cA_t(\varphi) = \langle \varphi, \int_0^t  F(u_s)  \dd s \rangle
$$
for any test function $\varphi \in \mathcal S$.
\end{definition}

In order to show these properties of controlled processes we will need some stochastic calculus. So let us recall here some basic elements needed below.

 
For any test function $\varphi\in\cS$ the processes $(u_t(\varphi))_{t}$ and $(\hat u_t(\varphi))_{t}$ are Dirichlet processes: sums of a martingale and a zero quadratic variation process. 
Note that we do not want to assume controlled processes to be semimartingales (even when tested with smooth functions). This is compatible with the regularity of our solutions and there is no clue that solutions of the SBE even with $\theta=1$ are distributional semimartingales. A suitable notion of stochastic calculus which is valid for a large class of processes and in particular for Dirichlet processes is the stochastic calculus via regularization developed by Russo and Vallois~\cite{RV}. In this approach the It\^o formula can be extended to Dirichlet processes. In particular if $(X^i)_{i=1,\dots,k}$ is an $\RR^k$ valued Dirichlet process and $g$ is a $C^2(\RR^k;\RR)$ function then
$$
g(X_t) = g(X_0) + \sum_{i=1}^k\int_0^t \partial_i g(X_s) \dd^- X^i_s + \frac12 \sum_{i,j=1}^k \int_0^t \partial^2_{i,j} g(X_s) \dd^- [X^i,X^j]_s
$$
where $\dd^-$ denotes the forward integral and $[X,X]$ the quadratic covariation of the vector process $X$. Decomposing $X=M+N$ as the sum of a martingale $M$ and a zero quadratic variation process $N$ we have $[X,X]=[M,M]$ and
$$
g(X_t) = g(X_0) +  \sum_{i=1}^k \int_0^t \partial_i g(X_s) \dd^- M^i_s +  \sum_{i=1}^k \int_0^t \partial_i g(X_s) \dd^- N^i_s
$$
$$ +   \sum_{i,j=1}^k\frac12 \int_0^t \partial^2_{i,j} g(X_s) \dd^- [M^i,M^j]_s
$$
where now $\dd^- M$ coincide with  the usual It\^o integral and $[M,M]$ is the usual quadratic variation of the martingale $M$. The integral $\int_0^t \partial_i g(X_s) \dd^- N^i_s$ is well-defined due to the fact that all the other terms in this formula are well defined. The case the function $g$ depends explicitly on time can be handled by the above formula by considering time as an additional (0-th) component of the process $X$ and using the fact that $[X^i,X^0]=0$ for all $i=1,..,k$. In the computations below we will only need to apply the It\^o formula to smooth functions.


\section{The It\^o trick}
\label{sec:ito-trick}
Let us denote by $L_0$ the generator of the Ornstein-Uhlenbeck process associated to the operator $A^\theta$:
\[
L_0 \varphi(x) = \sum_{k \in \mathbb Z_0} |k|^{2\theta} \big(- x_k D_k \varphi(x) + \tfrac{1}{2} D_{-k}D_k \varphi(x)\big). 
\]

Consider now $u\in\cR_\theta$ and  a smooth cylinder function  $h:[0,T]\times \Pi_N H\to \RR$. The It\^o formula for the finite quadratic variation process $u^N = \Pi_N u$ gives
$$
h(t,u^N_t)=h(0,u^N_0)+\int_0^t (\partial_s + L^N_0) h(s,u^N_s) \dd s +\int_0^t D h(s,u^N_s) \dd \cA_s + M^+_t 
$$ 
where 
$$
L^N_0 h(s,x) = \sum_{k\in\ZZ_0 : |k|\le N} |k|^{2\theta} ( x_{k} D_k h(s,x) + D_k D_{-k} h(s,x))
$$
is the restriction of the operator $L_0$ to $\Pi_N H$,
and the It\^o formula on the backward process gives
$$
h(T-t,u^N_{T-t})=h(T,u^N_T)+ \int_0^{t} (-\partial_s + L^N_0) h(T-s,u^N_{T-s}) \dd s
$$
$$
- \int_0^{t} D h(T-s,u^N_{T-s}) \dd \cA_{T-s} + M^-_t 
$$ 
so
$$
 -M^+_t + M^-_{T-t}-M^-_T = \int_0^t 2 L_0^N   h(s,u^N_{s})\dd s. 
$$


The martingale $M^+$ has quadratic variation given by
$
[ M^+ ]_t = \int_0^t \cE^\theta_N(h(s,\cdot))(u^N_s) \dd s
$,
where 
$$
\cE_N^\theta(\varphi)(x) = \frac12  \sum_{k\in \ZZ_0: |k|\le N}|k|^{2\theta}  |D_k \varphi(x)|^2 ,
$$
with a similar expression for $M^-$. Let
$$
\cE^\theta(\varphi)(x) = \frac12  \sum_{k\in \ZZ_0}|k|^{2\theta}  |D_k \varphi(x)|^2 ,
$$
then using Burkholder--Davis--Gundy inequality we can prove the following bound.
\begin{lemma}[It\^o trick] 
\label{lemma:ito-trick}
Let $h : [0,T]\times \Pi_N H \to \RR$ be a cylinder function. Then
for any $p \geq 1$,
\begin{equation}
\label{eq:ito-trick}
\left\|\sup_{t\in[0,T]}\left|\int_0^t L_0 h(s,\Pi_N u_s) \dd s\right|\right\|_{L^p(\PP_\mu)} \lesssim_p 
T^{1/2} \sup_{s\in[0,T]}\left\|  \cE^\theta(h(s,\cdot))  \right\|^{1/2}_{L^{p/2}(\mu)} .
\end{equation}
In the particular case $h(s,x)= e^{a(T- s)}\tilde h(x)$ for some $a\in\RR$  we have the improved estimate
\begin{equation}
\label{eq:ito-trick-conv}
\begin{split}
\left\|\int_0^T e^{a(T-s)} L_0 \tilde h(\Pi_N u_s) \dd s\right\|_{L^p(\PP_\mu)} 
\lesssim_p
 \left(\frac{1-e^{2aT}}{2a}\right)^{1/2} 
 \left\| \cE^\theta(\tilde h) \right\|^{1/2}_{L^{p/2}(\mu)}
 \end{split}
\end{equation}
\end{lemma}
\begin{proof}
$$
\left\|\sup_{t\in[0,T]}\left|\int_0^t  2 L_0^N h(s,u_s) \dd s\right|\right\|_{L^p(\PP_\mu)} \le
\left\|\sup_{t\in[0,T]}|M^+_t| \right\|_{L^p(\PP_\mu)}+
2 \left\|\sup_{t\in[0,T]}|M^-_t| \right\|_{L^p(\PP_\mu)}
$$
$$
\lesssim_p \left\| \langle M^+\rangle_T  \right\|_{L^{p/2}(\PP_\mu)}^{1/2}+\left\| \langle M^-\rangle_T  \right\|_{L^{p/2}(\PP_\mu)}^{1/2}
\lesssim_p \left\|\int_0^T \cE^\theta(h(s,\cdot))(u_s) \dd s \right\|_{L^{p/2}(\PP_\mu)}^{1/2}
$$
$$
\lesssim_p \left(\int_0^T\left\| \cE^\theta(h(s,\cdot))(u_s) \right\|_{L^{p/2}(\PP_\mu)} \dd s\right)^{1/2}
\lesssim_p T^{1/2} \sup_{s\in[0,T]} \left\| \cE^\theta(h(s,\cdot)) \right\|_{L^{p/2}(\mu)}^{1/2} .
$$
For the convolution we bound as follows
\begin{equation*}
%\label{eq:ito-trick-conv}
\begin{split}
\left\|\int_0^T e^{a(T-s)} 2 L_0^N \tilde h(u_s) \dd s\right\|_{L^p(\PP_\mu)} 
&  \lesssim_p \left(\int_0^T e^{2a(T-s)}\dd s\right)^{1/2} 
 \left\| \cE^\theta(\tilde h)(u_0) \right\|^{1/2}_{L^{p/2}(\PP_\mu)} 
 \\ &  \lesssim_p
 \left(\frac{1-e^{2aT}}{2a}\right)^{1/2} 
 \left\| \cE^\theta(\tilde h) \right\|^{1/2}_{L^{p/2}(\mu)}
 \end{split}
\end{equation*}
\end{proof}
The bound~\eqref{eq:ito-trick} in the present form (with the use of the backward martingale to remove the drift part) it has been inspired by~\cite{CLO}*{Lemma 4.4}.

\begin{lemma}[Exponential integrability] 
Let  $h : [0,T]\times \Pi_N H \to \RR$ be a cylinder function. Then
\begin{equation}
\label{eq:exp-ito-trick}
\expect \sup_{t\in[0,T]}e^{2 \int_0^t L_0^N h(s,\Pi_N u_s) \dd s} \lesssim 
\expect e^{8 \int_0^T \cE^\theta(h(s,u_s)) \dd s } 
\end{equation}
\end{lemma}
\begin{proof}
Let as above $M^\pm$ be the (Brownian) martingales in the representation of the integral $\int_0^t L_0^N h(s,\Pi_N u_s) \dd s$. By Cauchy-Schwartz
$$
\expect \sup_{t\in[0,T]}e^{2 \int_0^t L_0^N h(s,\Pi_N u_s) \dd s} \le \left[\expect \sup_{t\in[0,T]}e^{2 M^+_t}\right]^{1/2}  \left[\expect \sup_{t\in[0,T]}e^{2 (M^-_T-M^-_{T-t})}\right]^{1/2}.
$$
By  Novikov's criterion 
$
e^{4 M^+_t - 8 \langle M^+\rangle_t }
$ is a martingale for $t\in[0,T]$ if $\expect e^{8 \langle M^+\rangle_T} < \infty$. In this case
$$
\expect \sup_{t\in[0,T]}e^{2 M^+_t} \le \expect \sup_{t\in[0,T]}(e^{2 M^+_t- 4 \langle M^+\rangle_t}\sup_{t\in[0,T]} e^{ 4 \langle M^+\rangle_t })
$$
$$
\le\left[\expect \sup_{t\in[0,T]} e^{4 M^+_t- 8 \langle M^+\rangle_t}\right]^{1/2} \left[\expect e^{8 \langle M^+\rangle_T }\right]^{1/2}
$$
and by Doob's inequality we get that the previous expression is bounded by
$$
\left[\expect e^{4 M^+_T- 8 \langle M^+\rangle_T}\right]^{1/2} \left[\expect e^{8 \langle M^+\rangle_T }\right]^{1/2}
\le \left[\expect e^{8 \langle M^+\rangle_T }\right]^{1/2}.
$$
Reasoning similarly for $M^-$ we obtain that 
$$
\expect \sup_{t\in[0,T]}e^{2 \int_0^t L_0^N h(s,\Pi_N u_s) \dd s} \le \expect e^{8 \langle M^+\rangle_T } = \expect e^{8 \int_0^T \cE^\theta(h(s,u_s)) \dd s }.
$$
\end{proof}

\section{Estimates on the Burgers drift}
\label{sec:estimates}

In this section we provide the key estimates on the Burgers drift via the quadratic variations of the forward and backward martingales in its decomposition.
Let $F(x)(\xi) = B (x(\xi))^2$. Define 
$$
H(x) = -\int_0^\infty F(e^{-A^\theta t}x)\dd t
$$
and consider $L_0 H(x)$ as acting on each Fourier coordinate of $H(x)$. Then it is easy to check that 
$$
L_0 H(\Pi_N x) = \langle A^\theta x, D H(\Pi_N x)\rangle  = -2 \int_0^\infty B [(e^{-A^\theta t}\Pi_N x)(A^\theta e^{-A^\theta t} \Pi_N x) ]\dd t
$$
$$
=  -\int_0^\infty \frac{\dd }{\dd t}B [(e^{-A^\theta t}\Pi_N x)^2 ] = B (\Pi_N x)^2=F(\Pi_N x)
$$
since $\lim_{t\to \infty} B [(e^{-A^\theta t}\Pi_N x)^2 ]  = 0$. Let $H_N(x)= H(\Pi_N x)$ then this function in Fourier coordinates reads
$$
(H_{N}(x))_k =
2 ik  \sum_{k_1,k_2 : k=k_1+k_2}  \frac{\II_{|k|,|k_1|,|k_2|\le N}}{|k_1|^{2\theta}+|k_2|^{2\theta}} x_{k_1} x_{k_2}.
$$
Let us denote with $(H_{N}(x))_k^{\pm}$ respectively the real and imaginary parts of this quantity: $(H_{N}(x))_k^{\pm}= ((H_{N}(x))_k\pm (H_{N}(x))_{-k})/(2 i^{\pm})$ where $i^+=1$ and $i^-=i$. Now
$$
(H_{N}(x))^\pm_k =
 i^{\mp}k  \sum_{k_1,k_2 : k=k_1+k_2}  \frac{\II_{|k|,|k_1|,|k_2|\le N}}{k_1^{2\theta}+k_2^{2\theta}} (x_{k_1} x_{k_2}\mp x_{-k_1} x_{-k_2})
$$
and recall that
$
\cE^\theta((H_N)^\pm_k)(x) = \sum_{q\in\ZZ_0} |q|^{2\theta} |D_q H^\pm_{N,k}(x)|^2
$.
\begin{lemma}
\label{lemma:energy-estimates}
For $\lambda >0$ small enough we have
\begin{equation}
\sup_{k\in\ZZ_0} \expect \exp\left[\lambda |k|^{2\theta-3}  \cE^\theta((H_N)^\pm_k)(u_0)\right] \lesssim 1
\end{equation}
and
\begin{equation}
\sup_{1\le M \le N}\sup_{k\in\ZZ_0} \expect \exp\left[\lambda |k|^{-2} M^{2\theta-1}  \cE^\theta((H_N-H_M)^\pm_k)(u_0)\right] \lesssim 1.
\end{equation}
\end{lemma}
\begin{proof}
We start by computing $\cE((H_N)^\pm_k)$: noting that
$$
D_q (H_{N})^\pm_k(x)=  i^\mp k  \left[  \frac{\II_{|k|,|q|,|k-q|\le N}}{|q|^{2\theta}+|k-q|^{2\theta}}  x_{k-q}\mp \frac{\II_{|k|,|q|,|k+q|\le N}}{|q|^{2\theta}+|k+q|^{2\theta}}  x_{-k-q}\right]
$$ 
we have
\begin{equation*}
%\label{eq:dirichlet}
\begin{split}
 \cE^\theta((H_N)^\pm_k)(x) & = \sum_{q\in\ZZ_0} |k|^2 |q|^{2\theta}\left[
 2 \frac{\II_{|k|,|q|,|k-q|\le N}}{(|q|^{2\theta}+|k-q|^{2\theta})^2}  |x_{k-q}|^2
 \right . \\
  & \left .
 \qquad \qquad \mp \frac{\II_{|k|,|q|,|k-q|\le N}}{|q|^{2\theta}+|k-q|^{2\theta}}  \frac{\II_{|k|,|q|,|k+q|\le N}}{|q|^{2\theta}+|k+q|^{2\theta}}  (x_{k-q} x_{k+q}+x_{-k+q} x_{-k-q})
 \right]
\end{split}
\end{equation*}
which gives the bound
$$
 \cE^\theta((H_N)^\pm_k)(x)\lesssim   |k|^2  \sum_{\substack{k_1,k_2 :k_1+k_2=k\\|k|,|k_1|,|k_2|\le N}} \frac{|k_1|^{2\theta}\II_{|k|,|k_1|,|k_2|\le N}}{(|k_1|^{2\theta}+|k_2|^{2\theta})^2}  |x_{k_2}|^2 
$$
$$
\lesssim  |k|^2  \sum_{\substack{k_1,k_2 :k_1+k_2=k\\|k|,|k_1|,|k_2|\le N}} \frac{\II_{|k|,|k_1|,|k_2|\le N}}{|k_1|^{2\theta}+|k_2|^{2\theta}}  |x_{k_2}|^2 =   \sum_{\substack{k_1,k_2 :k_1+k_2=k\\|k|,|k_1|,|k_2|\le N}}c(k,k_1,k_2) |x_{k_2}|^2  = h_N(x)
$$
where $c(k,k_1,k_2) = |k|^2/(|k_1|^{2\theta}+|k_2|^{2\theta})$.
Let $$
I_N(k) =  \sum_{\substack{k_1,k_2 :k_1+k_2=k\\|k|,|k_1|,|k_2|\le N}}  c(k,k_1,k_2) 
$$
and note that the sum in $I_N(k)$ can be bounded by the equivalent integral giving (uniformly in $N$)
$$
I_N(k) \lesssim  |k|^{2}  \int_{\RR} \frac{\dd q}{|q|^{2\theta}+|k-q|^{2\theta}} 
= |k|^{3-2\theta}  \int_{\RR} \frac{\dd q}{|q|^{2\theta}+|1-q|^{2\theta}} \lesssim |k|^{3-2\theta} 
$$
since that the last integral is finite for $\theta > 1/2$. Then
$$
\expect e^{\lambda |k|^{2\theta-3}\cE^\theta((H_N)^\pm_k)(u_0)} \le \expect e^{\lambda C|k|^{2\theta-3} h_N(u_0)}
$$
$$
 \le \sum_{\substack{k_1,k_2 :k_1+k_2=k\\|k|,|k_1|,|k_2|\le N}} c(k,k_1,k_2)  \expect \frac{e^{\lambda C |k|^{2\theta-3} I_N(k) |(u_0)_{k_2}|^2}}{I_N(k)}
\le \sum_{\substack{k_1,k_2 :k_1+k_2=k\\|k|,|k_1|,|k_2|\le N}} c(k,k_1,k_2)  \expect \frac{e^{\lambda C'|(u_0)_{k_2}|^2}}{I_N(k)} 
$$
where we used the previous bound to say that $C|k|^{2\theta-3} I_N(k) \le C'$ uniformly in $k$. Remind that $(u_0)_k$ has a Gaussian distribution of mean zero and unit variance. Therefore for $\lambda$ small enough $\expect e^{\lambda C'|(u_0)_{k_2}|^2}\lesssim 1$ uniformly in $k_2$ so that  
$$
\expect e^{\lambda |k|^{2\theta-3}\cE^\theta((H_N)^\pm_k)(u_0)} \lesssim 1.
$$
This establishes the claimed exponential bound for $\cE^\theta((H_N(x))_k^\pm)$. 
Similarly we have
$$
\cE^\theta((H_N-H_M)^\pm_k)(x)  \lesssim \sum_{k_1,k_2 :k_1+k_2=k} (\II_{|k|,|k_1|,|k_2|\le N}-\II_{|k|,|k_1|,|k_2|\le M})^2 c(k,k_1,k_2) | x_{k_2}|^2 .
 $$
Let
$$
I_{N,M}(k) =\sum_{k_1,k_2 :k_1+k_2=k} (\II_{|k|,|k_1|,|k_2|\le N}-\II_{|k|,|k_1|,|k_2|\le M})^2 c(k,k_1,k_2)
$$
and note that, for $N\ge M$,
$$
(\II_{|k|,|k_1|,|k_2|\le N}-\II_{|k|,|k_1|,|k_2|\le M}) \lesssim \II_{|k|,|k_1|,|k_2|\le N}(\II_{|k|> M}+\II_{|k_1|> M}+\II_{|k_2|> M}).
$$
Then, by estimating the sums with the corresponding integrals and after easy simplifications we remain with the following bound
$$
I_{N,M}(k)\lesssim |k|^{2} \II_{|k|> M} \int_{\RR} \frac{\dd q}{|q|^{2\theta}+|k-q|^{2\theta}} 
+ |k|^{2}  \int_{\RR} \frac{\II_{|q|> M}\dd q}{|q|^{2\theta}+|k-q|^{2\theta}} 
$$
The first integral in the r.h.s.~is easily handled by
$$
|k|^{2} \II_{|k|> M} \int_{\RR} \frac{\dd q}{|q|^{2\theta}+|k-q|^{2\theta}} \lesssim |k|^{3-2\theta} \II_{|k|> M} \lesssim |k|^2 M^{1-2\theta}
$$
since $\theta > 1/2$. For the second we have the analogous bound
$$
|k|^{2}  \int_{\RR} \frac{\II_{|q|> M}\dd q}{|q|^{2\theta}+|k-q|^{2\theta}} \lesssim 
|k|^{2}  \int_{\RR} \frac{\II_{|q|> M}\dd q}{|q|^{2\theta}} \lesssim |k|^2 M^{1-2\theta}
$$
which concludes the proof.
\end{proof}


The hypercontractivity of the measure $\mu$ allows to obtain a first trivial bound on the drift. Let 
$$
G^M_t = \int_0^t F_{M}(u_s) \dd  s,
$$
where $F_M(u_t) = F(\Pi_M u_t)$. Then
$$
\|\sup_{t\in[0,T]}\left| (G^M_t)_k \right| \|_{L^p(\PP_\mu)}
\lesssim 
\int_0^T \left\| (F_{M}(u_s))_k \right\|_{L^p(\PP_\mu)}\dd  s 
$$
$$\lesssim 
T \left\| (F_{M}(u_0))_k  \right\|_{L^p(\PP_\mu)}
\lesssim_p 
  |k| M T.
$$
Then using Lemma~\ref{lemma:ito-trick}, the estimates contained in Lemma~\ref{lemma:energy-estimates} and letting  $F_{N,M}(x) = F_N(x)-F_M(x)$ we are led to the next set of more refined estimates for the drift and his small scale contributions.
\begin{lemma} For any $M\le N$ we have
\begin{equation}
\label{eq:basic-est-1}
\|\sup_{t\in[0,T]} \left|(G^M_t)_k\right| \|_{L^p(\PP_\mu)}\lesssim_p |k| M T
\end{equation}
\begin{equation}
\label{eq:basic-est-2}
\|\sup_{t\in[0,T]} \left|(G^M_t)_k\right| \|_{L^p(\PP_\mu)}\lesssim_p |k|^{3/2-\theta} T^{1/2}
\end{equation}
\begin{equation}
\label{eq:basic-est-3}
\|\sup_{t\in[0,T]}\left|(G^M_t)_k-(G^N_t)_k\right| \|_{L^p(\PP_\mu)}\lesssim_p |k| T^{1/2} M^{1/2-\theta}
\end{equation}
\begin{equation}
\label{eq:basic-est-4}
\|\sup_{t\in[0,T]}\left|(G^M_t)_k\right| \|_{L^p(\PP_\mu)}\lesssim_p   |k| T^{2\theta/(1+2\theta)}
\end{equation}
\end{lemma}

Where the last bound is obtained from the previous two by decomposing $F_N(x) = F_M(x)+F_{N,M}(x)$ and performing the optimal choice $M \sim T^{-1/(1+2\theta)}$.

These estimates go throught also for the functions obtained via convolution with the $e^{-A^\theta t}$ semi-group using eq.~\eqref{eq:ito-trick-conv}. Let
$$
\tilde G^M_t = \int_0^t e^{-A^\theta (t-s)} F_{M}(u_s) \dd  s
$$
then
\begin{lemma} For any $M\le N$ we have
\begin{equation}
\label{eq:basic-est-1-conv}
\|(\tilde G^M_t)_k\|_{L^p(\PP_\mu)}\lesssim_p |k| M \left(\frac{1-e^{-2k^{2\theta} t/2}}{2k^{2\theta}}\right)
\end{equation}
\begin{equation}
\label{eq:basic-est-2-conv}
\|(\tilde G^M_t)_k\|_{L^p(\PP_\mu)}\lesssim_p |k|^{3/2-\theta} \left(\frac{1-e^{-2k^{2\theta} t/2}}{2k^{2\theta}}\right)^{1/2}
\end{equation}
\begin{equation}
\label{eq:basic-est-3-conv}
\|(\tilde G^M_t)_k-(\tilde G^N_t)_k\|_{L^p(\PP_\mu)}\lesssim_p |k| M^{1/2-\theta} \left(\frac{1-e^{-2k^{2\theta} t/2}}{2k^{2\theta}}\right)^{1/2}
\end{equation}
\end{lemma}
In particular, the second of these inequalities  gives
$$
\|(\tilde G^N_t)_k \|_{L^p(\PP_\mu)}\lesssim_p |k|^{3/2-2\theta} 
$$

To control the time regularity of the drift convoluton we consider $0\le s \le t$ and decompose
$$
\|(\tilde G^N_t)_k - (\tilde G^N_s)_k\|_{L^p(\PP_\mu)}
$$
$$
\le \| \int_s^t (e^{-A^\theta(t-r)} F_{N}(u_r))_k \dd  r\|_{L^p(\PP_\mu)} + (e^{-k^{2\theta}(t-s)}-1)\|(\tilde G^N_s)_k\|_{L^p(\PP_\mu)}
$$
$$
\lesssim |k|^{3/2-\theta} (t-s)^{1/2} +|k|^{3/2-2\theta}(e^{-k^{2\theta}(t-s)}-1)
\lesssim |k|^{3/2-\theta} (t-s)^{1/2}
$$
By interpolation with the estimate
$$
\|(\tilde G^N_s)_k - (\tilde G^N_s)_k\|_{L^p(\PP_\mu)} \lesssim_{p} |k|^{3/2-2\theta}
$$
we get 
\begin{lemma}
For all small $\eps > 0$
\begin{equation}
\label{eq:basic-est-4-conv}
\|(\tilde G^N_s)_k - (\tilde G^N_s)_k\|_{L^p(\PP_\mu)} \lesssim_{p} |k|^{3/2-2\theta+2\eps \theta} (t-s)^\eps
\end{equation}
\end{lemma}

All these $L^p$ estimates can be replaced with equivalent exponential estimates. For example it is not difficult to prove that for small $\lambda$ we have
$$
\sup_{t\in[0,T]} \sup_{k\in\ZZ_0} \expect \exp\left(\lambda |k|^{2\theta-3/2} (\tilde G^N_t)^\pm_k \right) \lesssim 1
$$
where $(\cdot)^\pm$ denote, as before, the real and imaginary parts, respectively.

\section{Existence of controlled solutions}
\label{sec:existence}



Consider the SDE on $H$ given by
\begin{equation}
\label{eq:burgers-reg}
\dd u^N_t = - A^\theta u^N_t  \dd t + F_N(u^N_s)\dd t + A^{\theta/2} \dd W_t,
\end{equation}
where  $F_N : H\to H$ is defined by $F_N(x) = \frac12  \Pi_N B (\Pi_N x)^2$. The evolution described by~\eqref{eq:burgers-reg} preserves the Gaussian measure $\mu$ on $H=\Lz$ with covariance $\mathrm{Id}$ which is characterized by the equation
$$
\int e^{i \langle \psi,x \rangle}\mu(\dd x) = e^{-\langle \psi,\psi\rangle/2}\qquad \psi\in H
$$
or alternatively the integration by parts formula
$$
\int D_k \varphi(x) \mu(\dd x) = \int x_{-k} \varphi(x) \mu(\dd x) 
$$
for all $k\in\ZZ_0$.
The diffusion $u^N$ has  generator
$$
L_N \varphi(x) = L_0\varphi(x)+  \sum_{k\in\ZZ_0, |k|\le N} (F_N(x))_k D_k \varphi(x)
$$
where
$$
L_0\varphi(x)= \sum_{k\in\ZZ_0} \left[\frac12 k^{2\theta} D_{-k}D_k \varphi(x) + (A^\theta x)_k D_k \varphi(x)\right] 
$$
is the generator of the Ornstein--Uhlenbeck (OU) process which satisfy the integration by parts formula
$
\mu [\varphi L_0 \varphi] = \mu[ \cE(\varphi,\varphi)]
$
The mild formulation~\cite{DZ} of eq.~\eqref{eq:burgers-reg} is
\begin{equation}
\label{eq:burgers-reg-mild}
u^N_t = e^{-A^\theta t} u_0 +   \int_0^t e^{-A^\theta (t-s)} F_N(u^N_s) \dd s + A^{\theta/2} \int_0^t e^{-A^\theta (t-s)}  \dd W_s
\end{equation}
 where the stochastic convolution in the r.h.s is defined as
$$
A^{\theta/2} \int_0^t e^{-A^\theta (t-s)}  \dd W_s = \sum_{k\in\ZZ_0} |k|^{\theta} e_k \int_0^t e^{-|k|^{2\theta} (t-s)} \dd \beta^k_s .
$$

Now let
$$
\cA_t^{N}=\int_0^t F_{N}(u^{N}_s) \dd s ,\qquad \tilde \cA_t^{N}=\int_0^t e^{-A^\theta(t-s)} F_{N}(u^{N}_s) \dd s . 
$$

\begin{lemma}
Let $\sigma=(3/2-2\theta)_+$.
The processes $\{(u^N,\cA^N,\tilde \cA^N,W)\}_N$ are tight in the space of continuous functions with values in $\cF L^{\infty,\sigma-\eps}\times \cF L^{\infty,3/2-\theta-\eps}\times \cF L^{\infty,3/2-2\theta-\eps}\times \cF L^{\infty,-\eps}$ for all small $\eps > 0$.
\end{lemma}

\begin{proof}
The estimate~\eqref{eq:basic-est-4-conv} in the previous section readily gives that for any small $\eps >0$ and sufficienly large $p$
$$
\expect_\mu\left[\sum_{k\in\ZZ_0} |k|^{-(3/2-2\theta+3\theta\eps) p} \left(|(\tilde \cA_t^{N}-\tilde \cA_s^{N})_k| \right)^p\right]
\lesssim_{p,\eps} \sum_{k\in\ZZ_0} |k|^{-\theta\eps p} |t-s|^{p \eps} \lesssim |t-s|^{p \eps}
$$
This estimates show that the family of processes $\{ \tilde \cA^{N} \}_{N}$ is tight in $C([0,T],\cF L^{\infty,\alpha})$ for $\alpha=3/2-2\theta+3\theta\eps$ and sufficiently small $\eps>0$.
An analogous argument using the estimate~\eqref{eq:basic-est-2}
 shows that the family of processes $\{  \cA^{N} \}_{N}$ is tight in $C^\gamma([0,T],\cF L^{\infty,\beta})$ for any $\gamma<1/2$ and $\beta < 3/2-\theta$.
It is not difficult to show that the stochastic convolution $\int_0^t e^{-A^\theta(t-s)} A^{\theta/2} \dd W_s$ belongs to $C([0,T],\cF L^{\infty,1-\theta-\eps})$ for all small  $\eps>0$. 
Taking into account the mild equation~\eqref{eq:burgers-reg-mild}
we find that the processes $\{(u^{N}_t)_{t\in[0,T]}\}_{N}$ are tight in $C([0,T],\cF L^{\infty,\sigma-\eps})$. 
\end{proof}


\begin{theorem}
Any accumulation point $(u,\cA,\tilde\cA,W)$ of the family $\{(u^N,\cA^N,\tilde\cA^N , W)\}_N$ has stationary $u$ and satisfy
\begin{equation}
\label{eq:limit-1}
u_t =  u_0 +  \cA_t  -  \int_0^t A^\theta u_s \dd s +  B W_t = u_0 + \tilde \cA_t  + \int_0^t e^{-A^\theta(t-s)} B \dd W_s 
\end{equation}
where, as space distributions,
\begin{equation}
\label{eq:limit-4}
\cA_t = \lim_{M \to \infty}\int_0^t F_{M}(u_s) \dd s
\quad \text{ and } \quad 
\tilde \cA_t  = \int_0^t e^{-A^\theta(t-s)} \dd \cA_s .
\end{equation}
where this last integral can be defined as a Young integral. Moreover the process $u$ is a controlled process.
\end{theorem}

\begin{proof}
Let us first prove~\eqref{eq:limit-4}. Note that there exists a subsequence $\{N_n\}_n$ for which $u = \lim_n u^{N_n}$, then
$$
\int_0^t F_{M}(u_s) \dd s = \int_0^t (F_{M}(u_s) - F_{M}(u^{N_n}_s)) \dd s 
$$
$$\qquad + \int_0^t (F_{M}(u^{N_n}_s) - F_{N_n}(u^{N_n}_s)) \dd s
+ \int_0^t F_{N_n}(u^{N_n}_s) \dd s
$$
but now, in  $\cF L^{\infty,3/2-\theta-\eps}$,
$$
\lim_n \int_0^t F_{N_n}(u^{N_n}_s) \dd s = \cA_t
\qquad \lim_n \int_0^t (F_{M}(u_s) - F_{M}(u^{N_n}_s)) \dd s = 0 ,
$$
where for the second limit note that the functionals depends only of a finite number of components of $u$ and $u^{N_n}$ and that we have the distributional convergence of $u^{N_n}$ to $u$.
Moreover, for all $k\in\ZZ_0$,
$$
\lim_M \sup_{N_n : M<N_n} \left\| \int_0^t (F_{M}(u^{N_n}_s) - F_{N_n}(u^{N_n}_s))_k \dd s\right\|_{L^p(\PP_\mu)}= 0.
$$
Then note that, by the apriori estimates, $\cA^{N_n}$ converges tp $\cA$ in $C^\gamma(\cF L^{\infty,3/2-\theta-\eps}$ for all $\gamma <1/2$ and $\eps > 0$ so that we can use Young integration to define $\int_0^t e^{-A^\theta(t-s)} \dd \cA^{N_n}_s$ as a space distribution and to obtain its convergence (for example for each of its Fourier components) to $\int_0^t e^{-A^\theta(t-s)} \dd \cA^{N_n}_s$. At this point eq.~\eqref{eq:limit-1} is a simple consequence. The backward processes $\hat u^{N_n}_{t}=u^{N_n}_{T-t}$ and $\hat \cA^{N_n}_t = -\cA^{N_n}_{T-t}$ converge to $\hat u_{t}=u_{T-t}$ and $\hat \cA_t = -\cA_{T-t}$ respectively and moreover note that $\cA$ as a distributional process has trajectories which are H\"older continuous for any exponent smaller than $2\theta/(1+2\theta)>1/2$ as a consequence of the estimate~\eqref{eq:basic-est-4} and this directly implies that $\cA$ has zero quadratic variation. So $u$ is a controlled process in the sense of our definition.
\end{proof} 


\section{Uniqueness for $\theta>5/4$}
\label{sec:uniq}

In this section we prove a simple uniqueness result for controlled solutions which is valid when $\theta > 5/4$.


\begin{theorem}
\label{th:uniqueness}
Controlled solutions are unique when $\theta > 5/4$.
\end{theorem}
\begin{proof}
Let $u$ be a controlled solution to the equation and let $u^N$ be the Galerkin
approximations defined above with respect to the Brownian motion obtained from the martingale part of the decomposition of $u$ as a controlled process. Then,  by bilinearity,
$$
 F_N \left( u \right) - F_N \left( u^N \right)  = F_N ( \Pi_N u_s+u^N_s,\Delta^N_s)
$$
and the difference $\Delta^N = \Pi_N \left( u - u^N \right)$
 satisfies the equation
$$
\Delta^N_t = \Pi_N \int_0^t e^{- A^{\theta} ( t - s )} 
   F_N ( u_s+u^N_s,\Delta^N_s) \dd s +\varphi^N_t
$$
where 
$$
\varphi^N_t = \int_0^t e^{- A^{\theta} \left( t - s \right)} \left( F_{} \left( u \right)
   - F_N \left( u^{} \right) \right) \dd s
$$
Note that
\[ \| \sup_{t \in [ 0, T ]} | ( \varphi^N_t )_k | \|_{L^p (
   \mathbb{P}_{\mu} )} \lesssim_p \max(| k |^{1 - 2 \theta} N^{1
   / 2 - \theta},| k |^{3/2 - 2 \theta}) \]
which by interpolation gives
\[ \| \sup_{t \in [ 0, T ]} | ( \varphi^N_t )_k | \|_{L^p (
   \mathbb{P}_{\mu} )}\lesssim_p | k |^{3/2 - 2 \theta+\varepsilon} N^{-\varepsilon} \]
for any small $\eps >0$. Let
$$
\Phi_N = \sup_{k\in\ZZ_0} \sup_{t \in [ 0, T ]} |k|^{2\theta-3/2-2\eps} | ( \varphi^N_t )_k | 
$$   
then
$$
\expect \sum_{N>1} N \Phi_N^p \le\sum_{N>1} N \sum_{k\in\ZZ_0}  \sup_{t \in [ 0, T ]} |k|^{p(2\theta-3/2-2\eps)}\expect  | ( \varphi^N_t )_k |^p 
$$
$$
\lesssim_p \sum_{N>1} N^{1-\eps p} \sum_{k\in\ZZ_0} |k|^{-p\eps}<+\infty
$$   
for $p$ large enough, which implies that almost surely
$
\Phi_N \lesssim_{p,\omega} N^{-1/p} 
$.
For the other term we have
$$
\sup_{t\in[0,T]}\left|\left( \int_0^t e^{- A^{\theta} \left( t - s \right)} F_N \left(
   \Pi_N u + u^N, \Delta_N \right) \dd s \right)_k\right| 
  \lesssim A_N |k|^{3/2-2\theta+2\eps} Q_T
$$   
where
   $
    A_N = \sup_{t \in \left[ 0, T \right]} \sup_k \left| k \right|^{2 \theta -
   3/2 - 2\varepsilon} \left| \left( \Delta^N_t \right)_k \right| $
and
$$
Q_T = \sup_{t\in[0,T]} |k|^{2\theta-1/2-2\eps} \int_0^t e^{- |k|^{2\theta} \left( t - s \right)} \sum_{q\in \ZZ_0} |(
   \Pi_N u_s + u^N_s)_q| |k-q|^{3/2-2\theta+2\eps}  \dd s 
 $$
This gives
\[ 
A_N \leqslant   Q_T A_N + \Phi_N.
 \] 
Since $3/2-2 \theta <-1 $ (that is $\theta > 5/4$), we have the estimate:
 $$
Q_T \lesssim  \sup_{t\in[0,T]} |k|^{2\theta-1/2-2\eps} \left[\int_0^t e^{- p' |k|^{2\theta} \left( t - s \right)}  \dd s  \right]^{1/p'} \left[ \int_0^T  \sum_{q\in \ZZ_0}\frac{ |(
   \Pi_N u_s + u^N_s)_q|^{p}}{ |k-q|^{-3/2+2\theta-2\eps}}  \dd s \right]^{1/p}
 $$
valid for some $p> 1$ (with $1/p'+1/p=1$). Then
$$ Q_T \lesssim   |k|^{2\theta-1/2-2\eps-2\theta/p'}  \left[ \int_0^T  \sum_{q\in \ZZ_0}\frac{ |(
   \Pi_N u_s + u^N_s)_q|^{p}}{ |k-q|^{-3/2+2\theta-2\eps}}  \dd s \right]^{1/p}
$$
and taking $p$ large enough such that $2\theta-1/2-2\eps-2\theta/p'\le 0$ we obtain
$$
Q_T \lesssim_p \left[ \int_0^T  \sum_{q\in \ZZ_0}\frac{ |(
   \Pi_N u_s + u^N_s)_q|^{p}}{ |k-q|^{-3/2+2\theta-2\eps}}  \dd s \right]^{1/p}
$$
By the stationarity of the processes $u$ and $u^N$ and the fact that their marginal laws are the white  noise we have
$$
\expect[ Q_T^p] \lesssim_p \int_0^T  \sum_{q\in \ZZ_0}\frac{\expect |(
   \Pi_N u_s + u^N_s)_q|^{p}}{ |k-q|^{-3/2+2\theta-2\eps}}  \dd s = T \sum_{q\in \ZZ_0}\frac{1}{ |k-q|^{-3/2+2\theta-2\eps}} \lesssim_p T
$$
Then by a simple Borel-Cantelli argument, almost surely $Q_{1/n} \lesssim_{p,\omega} n^{-1+1/p}$. Putting together the estimates for $\Phi_N$ and that for $Q_{1/n}$ we see that there exists a (random)  $T$ such that $C Q_T\le 1/2$ almost surely and that for this $T$:
$
A_N \leqslant  2 \Phi_N
$, 
which given the estimate on $\Phi_N$ implies that $A_N \to 0$ as $N\to\infty$ almost surely and that the solution of the equation is unique and is the
(almost-sure) limit of the Galerkin approximations.
\end{proof}


\section{Alternative equations}
\label{sec:alternative}
The technique of the present paper extends straighforwardly to some other modifications of the stochastic Burgers equation.

\subsection{Regularization of the convective term}
 Consider for example the equation
\begin{equation}
\label{eq:burgers-daprato}
\dd u_t = - A u_t \dd t + A^{-\sigma}F(A^{-\sigma} u_t) \dd t + B \dd W_t 
\end{equation}
which is the equation considered by Da~Prato, Debbussche and Tubaro in~\cite{DDT}. Letting $F_\sigma(x) = A^{-\sigma}F(A^{-\sigma} x)$, denoting by $H_\sigma$ the corresponding solution of the Poisson equation and following the same strategy as above we obtain the same bounds
$$
 \cE((H_{\sigma,N})^\pm_k)(x)\lesssim     \sum_{\substack{k_1,k_2 :k_1+k_2=k\\|k|,|k_1|,|k_2|\le N}}c_\sigma(k,k_1,k_2) |x_{k_2}|^2  
$$
where $c_\sigma(k,k_1,k_2) = |k|^{2-4\sigma}/[|k_1|^{4\sigma}|k_1|^{4\sigma}(|k_1|^{2}+|k_2|^{2})]$.
This quantity can then be bounded in terms of the sum
$$
I_{\sigma,N}(k) = \sum_{\substack{k_1,k_2 :k_1+k_2=k\\|k|,|k_1|,|k_2|\le N}}c_\sigma(k,k_1,k_2) \lesssim |k|^{1-12\sigma}
$$
From which we can reobtain similar bounds to those exploited above.
For example
$$
\left\|\int_0^t (e^{-A(t-s)}F_{\sigma,M}(u_s))_k \dd  s \right\|_{L^p(\PP_\mu)}\lesssim_p |k|^{-1/2-6\sigma} 
$$
And in particular we have existence of controlled solutions when $8\sigma+2>1$, that is $\sigma>-1/8$ and uniqueness when $-1/2-6\sigma<-1$ that is $\sigma> 1/12$. Which is an improvement over the result in~\cite{DDT} which has uniqueness for $\sigma>1/8$. 

\subsection{The Sasamoto--Spohn discrete model}
Another application of the above techniques is to the analysis of the discrete approximation to the stochastic Burgers equation proposed by Spohn and Sasamoto in~\cite{Spohn}. Their model is the following:
\begin{equation}
\label{eq:sasamoto-spohn}
\begin{split}
\dd u_j & = (2N+1) (u_j^2+u_j u_{j+1}-u_{j-1}u_j-u^2_{j-1})\dd t
\\
 & \qquad  + (2N+1)^2(u_{j+1}-2 u_j+u_{j-1})\dd t + (2N+1)^{3/2} (\dd B_j - \dd B_{j-1})
\end{split}
\end{equation}
for $j=1,\dots,2N+1$ with periodic boundary conditions $u_0=u_{2N+1}$ and where the processes $(B_j)_{j=1,\dots,2N+1}$ are a family of independents standard Brownian motions with $B_0=B_{2N+1}$. This model has to be tought as the discretization of the dynamic of the periodic velocity field $u(x)$ with $x\in(-\pi,\pi]$ sampled on a grid of mesh size $1/(2N+1)$,  that is $u_j = u(\xi^N_j)$ with $\xi^N_j = -\pi+2\pi(j/(2N+1))$. This fixes also the scaling factors for the different contributions to the dynamics if we want that, at least formally, this equation goes to a limit described by a SBE. Passing to Fourier variables $\hat u(k) = (2N+1)^{-1}\sum_{j=0}^{2N-1} e^{i \xi^N_j k} u_j$ for $k\in \ZZ^N$ with $\ZZ^N = \ZZ\cap [-N,N]$ and imposing that $\hat u(0)=0$, that is, considering the evolution only with zero mean velocity we get the system of ODEs:
$$
\dd \hat u_t(k) = F^\flat_N(\hat u_t)_k \dd t 
 -  |g_N(k)|^2  \hat u_t(k) \dd t + (2N+1)^{1/2} g_N(k) \dd \hat B_t(k)
$$
for $k\in \ZZ_0^N=\ZZ_0\cap [-N,N]$, where $g_N(k)=(2N+1)(1-e^{i k/(2N+1)})$,
$$
F^\flat_N(u_t)_k = \sum_{k_1,k_2\in\ZZ^N_0} \hat u_t(k_1)\hat u_t(k_2)[ g_N(k)-g_N(k)^*+g_N(k_1)-g_N(k_2)^*]
$$
and $(\hat B_\cdot(k))_{k\in\ZZ_0^N}$ is a family of centred complex Brownian motions such that $\hat B(k)^* = \hat B(-k)$ and with covariance $\expect \hat B_t(k) \hat B_t(-l) = \II_{k=l} t (2N+1)^{-1}$. 
If we then let $\beta(k) = (2N+1)^{1/2} \hat B(k)$ we obtain a family of complex BM of covariance 
$\expect  \beta_t(k)  \beta_t(-l) = t \II_{k=l} $. The generator $L^\flat_N$ of this stochastic dynamics is given by
$$
L^{\flat}_N \varphi( x) = \sum_{k\in\ZZ^N_0} F^\flat_N(x)_k  D_k \varphi( x)+L^{g_N,OU}_N \varphi (x)
$$
with $$ 
L^{g_N}_N \varphi( x) = \sum_{k\in\ZZ^N_0} |g_N(xk)|^2 (x_k D_{k}+ D_{-k}  D_k) \varphi( x)
$$
the generator of the OU process corresponding to the linear part associated with the multiplier $g_N$. It is easy to check that the complete dynamics preserves the (discrete) white noise measure, indeed
$$
\sum_{k\in\ZZ_0^N} x_{-k} F^\flat_N(x)_k = \sum_{\substack{k,k_1,k_2\in\ZZ_0^N\\k+k_1+k_2=0}}  x_{k}  x_{k_1} x_{k_2}[ g_N(k)^*-g_N(k)+g_N(k_1)-g_N(k_2)^*] =0
$$
since the symmetrization of the r.h.s. with respect to the permutations of the variables $k,k_1,k_2$ yields zero. Then defining suitable controlled process with respect to the linear part of this equation we can prove our apriori estimates on additive functionals which are now controlled by the quantity
$$
 \cE^{g_N}((H_{g_N,N})^\pm_k)(x)\lesssim     \sum_{\substack{k_1,k_2 :k_1+k_2=k\\|k|,|k_1|,|k_2|\le N}}c_{g_N}(k,k_1,k_2) |x_{k_2}|^2  
$$
with $c_{g_N}(k,k_1,k_2) = |g_N(k)|^2/[(|g_N(k_1)|^{2}+|g_N(k_2)|^{2})]$. Moreover noting that
$$
|g_N(k)|^2 = 2 (2N+1)^2(1-\cos(2\pi k/(2N+1)) \sim |k|^2
$$
uniformly  $N$, it is possible to estimate this energy in the same way we did before in the case $\theta=1$ and obtain  that the family of stationary solutions of equation~\eqref{eq:sasamoto-spohn} is tight in $C([0,T],\cF L^{\infty,-\eps})$ for all $\eps >0$. Moreover using the fact that $g_N(k) \to ik$ as $N\to \infty$ uniformly for bounded $k$ and that 
$$
\pi_M F^\flat_N(\pi_M x)_k = \sum_{k_1,k_2\in\ZZ^N_0} \II_{|k|,|k_1|,|k_2|\le M}  x_{k_1} x_{k_2}[ g_N(k)-g_N(k)^*+g_N(k_1)-g_N(k_2)^*]
$$
$$
\to 3 i k \sum_{k_1,k_2\in\ZZ_0} \II_{|k|,|k_1|,|k_2|\le M}  x_{k_1} x_{k_2} = 3 F_M(x)_k
$$
it is easy to check that any accumulation point is a controlled solution of the stochastic Burgers equations~\eqref{eq:burgers-theta}. 

\section{2d stochastic Navier-Stokes equation}
\label{sec:ns}

We consider the problem of stationary solutions to the 2d stochastic Navier-Stokes equation considered in~\cite{albeverio-cruzeiro} (see also~\cite{albeverio-ferrario}). We would like to deal with invariant measures obtained by formally taking the kinetic energy of the fluid and considering the associated Gibbs measure. However this measure is quite singular and we need a bit of hyperviscosity in the equation to make our estimates work.
 
\subsection{The setting}
Fix $\sigma>0$ and consider the following stochastic differential equation
\begin{equation}
\label{eq:ns}
\dd (u_{t})_k = - |k|^{2+2\sigma} (u_{t})_k \dd t +  B_k(u_t) \dd t + |k|^{\sigma} \dd \beta^k_t
\end{equation}
where $(\beta^k)_{k\in \Q}$ is a family of complex BMs for which $(\beta^k)^* = \beta^{-k}$ and $\expect[\beta^k \beta^q] = \II_{q+k=0}$, $u$ is a stochastic process with continuous trajectories in the space of distributions on the two dimensional torus $\TT^2$,
$$
 B_k(x) = \sum_{k_1+k_2=k} b(k,k_1,k_2) x_{k_1} x_{k_2}
$$
where $x: \Q \to \CC$ is such that $x_{-k} = x_k^*$ and
$$
b(k,k_1,k_2) = \frac{(k^\bot \cdot k_1)(k \cdot k_2)}{k^2}
$$
with $(\xi,\eta)^\bot = (\eta,-\xi) \in \RR^2$. Apart from the two-dimensional setting and the difference covariance structure of the linear part this problem has the same structure as the one dimensional stochastic Burgers equation we considered before. 
Note that to make sense of it (and in order to construct controlled solutions)
we can consider the  Galerkin approximations constructed as follows. Fix $N$ and solve the  problem finite dimensional problem
\begin{equation}
\label{eq:ns-N}
\dd (u^N_{t})_k = - |k|^{2+\sigma} (u^N_{t})_k \dd t +  B^N_k(u^N_t) \dd t + |k|^{-\sigma} \dd \beta^k_t
\end{equation}
for $k \in \ZZ^2_N = \{k \in \ZZ^2 : |k| \le N\}$, where
\begin{equation}
\label{eq:drift-N}
 B^N_k(x) =\II_{|k|\le N} \sum_{\substack{k_1+k_2=k\\|k_1|\le N, |k_2|\le N}} b(k,k_1,k_2) x_{k_1} x_{k_2} 
 \end{equation}
The generator of the process $u^N$ is given by
$
L^N\varphi(x) = L_0\varphi(x) + \sum_{k\in \Q} B^N_k(x) D_k \varphi(x)
$
where
$$
L^0 \varphi(x) = \frac12  \sum_{k\in \Q} |k|^{2\sigma}( D_{-k}D_k\varphi(x) -|k|^2 x_k D_k \varphi(x))  
$$
is the generator of a suitable OU flow. Note moreover that the kinetic energy of $u$ given by
$
E(x) = \sum_k |k|^{2} |x_k|^2 
$
is invariant under the flow generated by $B^N$. Moreover
$
D_{k} B^N_k(x) = 0 
$
since $x_k$ does not enter in the expression of $B^N_k(x)$, so the vectorfields $B^N$ leave also the measure
$
\prod_{k \in \QN}  dx_k 
$
invariant. Then  the  (complex) Gaussian measures
$$
\gamma(dx) = \prod_{k \in \Q} Z_k e^{-|k|^2 |x_k|^2} dx_k
$$
is invariant under  the flow generated by $B^N$. (This measure should be understood restricted to the set $\{x \in \CC^{\Q} : x_{-k} = \overline{x_k} \}$).
The measure $\gamma$ is also invariant for the $u^N$ diffusion since it is invariant for $B^N$ and for the OU process generated by $L^0$.
Intoduce standard Sobolev norms
$
\|x\|_\sigma^2 = \sum_{k \in \Q} |k|^{2\sigma} |x_k|^2
$
and denote with $H^\sigma$ the space of elements $x$ with $\|x \|_{\sigma}<\infty$.
The measure $\gamma$ is the Gaussian measure associated to $H^1$ and is supported on any $H^\sigma$ with $\sigma<0$
$$
\int \|x\|_\sigma^2 \gamma(dx) = \sum_{k \in \Q}  |k|^{2\sigma-2} < \infty
$$
so $(\gamma, H^1, \cap_{\eps < 0}H^\eps)$ is an abstract Wiener space in the sense of Gross. Note that the vectorfield  $B_k(x)$ in not defined on the support of $\gamma$. To give sense of controlled solutions to this equation we need to control
$$
 \cE((H_{N})^\pm_k)(x)\lesssim     \sum_{\substack{k_1,k_2 :k_1+k_2=k\\|k|,|k_1|,|k_2|\le N}}c_{\text{ns}}(k,k_1,k_2) |x_{k_2}|^2  
$$
with $c_{\text{ns}}(k,k_1,k_2) = |k_1|^{2\sigma} |k_1|^2|k_2|^2/(|k_1|^{2+2\sigma}+|k_2|^{2+2\sigma})^2$ and note that the stationary expectation of this term can be estimated by
$$
I_N(k) =  \sum_{\substack{k_1,k_2 :k_1+k_2=k\\|k|,|k_1|,|k_2|\le N}}c_{\text{ns}}(k,k_1,k_2) |k_2|^{-2} \lesssim \sum_{\substack{k_1,k_2 :k_1+k_2=k\\|k|,|k_1|,|k_2|\le N}}\frac{|k_1|^{2+2\sigma}}{ (|k_1|^{2+2\sigma}+|k_2|^{2+2\sigma})^2}\lesssim 
$$
$$
 \lesssim \sum_{\substack{k_1,k_2 :k_1+k_2=k\\|k|,|k_1|,|k_2|\le N}}\frac{1}{ |k_1|^{2+2\sigma}+|k_2|^{2+2\sigma}}\lesssim |k|^{-2\sigma}
$$
for any $\sigma > 0$. This estimate allows to apply our machinery and obtain stationary controlled solutions to this equation. 




%%%%%%%%%%%%%%%%%%%%%%%%%%%%%%%%%%%%%%%%%%%%%%%%%%%%%%%%%%%%%%%%%%%%%%%%%%%%%%%%%%%%%%%%%%%%%%%%%%%%%%%%%%%%%%%%%%%%%%%%%%%%%%%
%%%%%%%%%%%%%%%%%%%%%%%%%%%%%%%%%%%%%%%%%%%%%%%%%%%%%%%%%%%%%%%%%%%%%%%%%%%%%%%%%%%%%%%%%%%%%%%%%%%%%%%%%%%%%%%%%%%%%%%%%%%%%%%
%%%%%%%%%%%%%%%%%%%%%%%%%%%%%%%%%%%%%%%%%%%%%%%%%%%%%%%%%%%%%%%%%%%%%%%%%%%%%%%%%%%%%%%%%%%%%%%%%%%%%%%%%%%%%%%%%%%%%%%%%%%%%%%
%%%%%%%%%%%%%%%%%%%%%%%%%%%%%%%%%%%%%%%%%%%%%%%%%%%%%%%%%%%%%%%%%%%%%%%%%%%%%%%%%%%%%%%%%%%%%%%%%%%%%%%%%%%%%%%%%%%%%%%%%%%%%%%
%%%%%%%%%%%%%%%%%%%%%%%%%%%%%%%%%%%%%%%%%%%%%%%%%%%%%%%%%%%%%%%%%%%%%%%%%%%%%%%%%%%%%%%%%%%%%%%%%%%%%%%%%%%%%%%%%%%%%%%%%%%%%%%

\begin{bibdiv} 
\begin{biblist}

\bib{albeverio-cruzeiro}{article}{
	title = {Global flows with invariant {(Gibbs)} measures for Euler and {Navier-Stokes} two dimensional fluids},
	volume = {129},
	issn = {0010-3616},
	url = {http://www.springerlink.com/content/u1406006h7x32575/abstract/},
	doi = {10.1007/BF02097100},
	number = {3},
	journal = {Communications in Mathematical Physics},
	author = {Albeverio, Sergio},
	author = {Cruzeiro, {Ana-Bela}},
	year = {1990},
	keywords = {Physics and Astronomy},
	pages = {431--444},
}

\bib{albeverio-ferrario}{incollection}{
	series = {Lecture Notes in Mathematics},
	title = {Some Methods of Infinite Dimensional Analysis in Hydrodynamics: An Introduction},
	volume = {1942},
	isbn = {978-3-540-78492-0},
	shorttitle = {Some Methods of Infinite Dimensional Analysis in Hydrodynamics},
	url = {http://www.springerlink.com/content/2n14280t40q7v34q/abstract/},
	booktitle = {{SPDE} in Hydrodynamic: Recent Progress and Prospects},
	publisher = {Springer Berlin / Heidelberg},
	author = {Albeverio, Sergio},
	author = {Ferrario, Benedetta},
	year = {2008},
	keywords = {Mathematics and Statistics},
	pages = {1--50},
}
\bib{assing1}{article} {
	title = {A Pregenerator for Burgers  Equation Forced by Conservative Noise},
	volume = {225},
	issn = {0010-3616},
	url = {http://www.springerlink.com/content/t48f9yxafdddjnwx/abstract/},
	doi = {10.1007/s002200100606},
	number = {3},
	journal = {Communications in Mathematical Physics},
	author = {Assing, Sigurd},
	year = {2002},
	keywords = {Physics and Astronomy},
	pages = {611--632},
}


\bib{assing2}{article} {
	title = {A rigorous equation for the {Cole-Hopf} solution of the conservative {KPZ} dynamics},
	url = {http://arxiv.org/abs/1109.2886},
	journal = {{arXiv:1109.2886}},
	author = {Assing, Sigurd},
	month = {sep},
	year = {2011},
}


\bib{babin-kdv}{article}{
	title = {On the regularization mechanism for the periodic Korteweg-de Vries equation},
	volume = {64},
	issn = {0010-3640},
	url = {http://dx.doi.org/10.1002/cpa.20356},
	doi = {10.1002/cpa.20356},
	number = {5},
	journal = {Communications on Pure and Applied Mathematics},
	author = {Babin, Anatoli V. and Ilyin, Alexei A. and Titi, Edriss S.},
	year = {2011},
	pages = {591?648}
}


\bib{babin-ns}{article}{
	title = {Regularity and integrability of {\$3\$D} Euler and {Navier-Stokes} equations for rotating fluids},
	volume = {15},
	issn = {0921-7134},
	number = {2},
	journal = {Asymptotic Analysis},
	author = {Babin, A. and Mahalov, A. and Nicolaenko, B.},
	year = {1997},
	pages = {103?150}
}


		\bib{BG}{article}{
   author={Bertini, L.},
   author={Giacomin, G.},
   title={Stochastic Burgers and KPZ equations from particle systems},
   journal={Comm. Math. Phys.},
   volume={183},
   date={1997},
   number={3},
   pages={571--607},
   issn={0010-3616},
   review={\MR{1462228 (99e:60212)}},
   doi={10.1007/s002200050044},
}
 
\bib{CLO}{article}{
   author={Chang, C.-C.},
   author={Landim, C.},
   author={Olla, S.},
   title={Equilibrium fluctuations of asymmetric simple exclusion processes
   in dimension $d\geq 3$},
   journal={Probab. Theory Related Fields},
   volume={119},
   date={2001},
   number={3},
   pages={381--409},
   issn={0178-8051},
   review={\MR{1821140 (2002e:60157)}},
   doi={10.1007/PL00008764},
}
		
\bib{DDT}{article}{
   author={Da Prato, Giuseppe},
   author={Debussche, Arnaud},
   author={Tubaro, Luciano},
   title={A modified Kardar-Parisi-Zhang model},
   journal={Electron. Comm. Probab.},
   volume={12},
   date={2007},
   pages={442--453 (electronic)},
   issn={1083-589X},
   review={\MR{2365646 (2008k:60147)}},
}

\bib{DF}{article}{
   author={Da Prato, G.},
   author={Flandoli, F.},
   title={Pathwise uniqueness for a class of SDE in Hilbert spaces and
   applications},
   journal={J. Funct. Anal.},
   volume={259},
   date={2010},
   number={1},
   pages={243--267},
   issn={0022-1236},
   review={\MR{2610386}},
   doi={10.1016/j.jfa.2009.11.019},
}
	
\bib{DZ}{book}{
   author={Da Prato, Giuseppe},
   author={Zabczyk, Jerzy},
   title={Stochastic equations in infinite dimensions},
   series={Encyclopedia of Mathematics and its Applications},
   volume={44},
   publisher={Cambridge University Press},
   place={Cambridge},
   date={1992},
   pages={xviii+454},
   isbn={0-521-38529-6},
   review={\MR{1207136 (95g:60073)}},
   doi={10.1017/CBO9780511666223},
}
		
\bib{FGP}{article}{
   author={Flandoli, F.},
   author={Gubinelli, M.},
   author={Priola, E.},
   title={Well-posedness of the transport equation by stochastic
   perturbation},
   journal={Invent. Math.},
   volume={180},
   date={2010},
   number={1},
   pages={1--53},
   issn={0020-9910},
   review={\MR{2593276}},
   doi={10.1007/s00222-009-0224-4},
}
	
\bib{MR1988703}{article}{
   author={Flandoli, Franco},
   author={Russo, Francesco},
   author={Wolf, Jochen},
   title={Some SDEs with distributional drift. I. General calculus},
   journal={Osaka J. Math.},
   volume={40},
   date={2003},
   number={2},
   pages={493--542},
   issn={0030-6126},
   review={\MR{1988703 (2004e:60110)}},
}

\bib{MR2065168}{article}{
   author={Flandoli, Franco},
   author={Russo, Francesco},
   author={Wolf, Jochen},
   title={Some SDEs with distributional drift. II. Lyons-Zheng structure,
   It\^o's formula and semimartingale characterization},
   journal={Random Oper. Stochastic Equations},
   volume={12},
   date={2004},
   number={2},
   pages={145--184},
   issn={0926-6364},
   review={\MR{2065168 (2006a:60105)}},
   doi={10.1163/156939704323074700},
}
		
	
\bib{JG}{article}{
  author = {{Goncalves}, P.},
  author= { {Jara}, M.},
    title = {Universality of KPZ equation},
  journal = {ArXiv e-prints},
  archivePrefix = {arXiv},
  eprint = {arXiv:1003.4478},
  primaryClass = {math.PR},
     year = {2010}
}
		
		
\bib{controlling}{article}{
	title = {Controlling rough paths},
	volume = {216},
	issn = {0022-1236},
	url = {http://dx.doi.org/10.1016/j.jfa.2004.01.002},
	doi = {10.1016/j.jfa.2004.01.002},
	number = {1},
	journal = {Journal of Functional Analysis},
	author = {Gubinelli, M.},
	year = {2004},
	pages = {86?140}
}
\bib{kdv}{article}{
	title = {Rough solutions for the periodic Korteweg?de Vries equation},
	volume = {11},
	issn = {1534-0392},
	number = {2},
	journal = {Communications on Pure and Applied Analysis},
	author = {Gubinelli, Massimiliano},
	year = {2012},
	pages = {709?733}
}

\bib{Hairer}{article}{
	title = {Solving the {KPZ} equation},
	url = {http://fr.arxiv.org/pdf/1109.6811v2},
	journal = {{ArXiv} e-prints},
	author = {Hairer, M.},
	month = {sep},
	year = {2011}
}
		
\bib{KPZ}{article}{
  title = {Dynamic Scaling of Growing Interfaces},
  author = {Kardar, Mehran  and Parisi, Giorgio  and Zhang, Yi-Cheng },
  journal = {Phys. Rev. Lett.},
  volume = {56},
  number = {9},
  pages = {889--892},
  numpages = {3},
  year = {1986},
  month = {Mar},
  doi = {10.1103/PhysRevLett.56.889},
  publisher = {American Physical Society}
}
		
\bib{KV}{article}{
   author={Kipnis, C.},
   author={Varadhan, S. R. S.},
   title={Central limit theorem for additive functionals of reversible
   Markov processes and applications to simple exclusions},
   journal={Comm. Math. Phys.},
   volume={104},
   date={1986},
   number={1},
   pages={1--19},
   issn={0010-3616},
   review={\MR{834478 (87i:60038)}},
}

\bib{MR2353387}{article}{
   author={Russo, Francesco},
   author={Trutnau, Gerald},
   title={Some parabolic PDEs whose drift is an irregular random noise in
   space},
   journal={Ann. Probab.},
   volume={35},
   date={2007},
   number={6},
   pages={2213--2262},
   issn={0091-1798},
   review={\MR{2353387 (2008j:60153)}},
   doi={10.1214/009117906000001178},
}
		
\bib{RV}{incollection}{
	series = {Lecture Notes in Mathematics},
	title = {Elements of Stochastic Calculus via Regularization},
	volume = {1899},
	isbn = {978-3-540-71188-9},
	url = {http://www.springerlink.com/content/63u35k7n446q1t64/abstract/},
	booktitle = {S�minaire de Probabilit�s {XL}},
	publisher = {Springer Berlin / Heidelberg},
	author = {Russo, Francesco and Vallois, Pierre},
	editor = {{Donati-Martin}, Catherine and �mery, Michel and Rouault, Alain and Stricker, Christophe},
	year = {2007},
	keywords = {Mathematics and Statistics},
	pages = {147--185}
}


\bib{Spohn}{article}{
	title = {Superdiffusivity of the {1D} Lattice {Kardar-Parisi-Zhang} Equation},
	volume = {137},
	issn = {0022-4715},
	url = {http://www.springerlink.com/content/54h7q2682842g701/abstract/},
	doi = {10.1007/s10955-009-9831-0},
	number = {5},
	journal = {Journal of Statistical Physics},
	author = {Sasamoto, Tomohiro and Spohn, Herbert},
	year = {2009},
	keywords = {Physics and Astronomy},
	pages = {917--935},
}
		
\end{biblist} 
\end{bibdiv} 

\end{document}